\chapter{Positive Jacobi fields and graphical convergence}\label{sec:positive-jacobi}
The stability inequality, the linearized mean-curvature equation, and the separation of nearby sheets are different statements. This section develops their connections. We use the scalar convention $H=\operatorname{div}_\Sigma\nu$ from Definition \ref{def:smooth-minimal}.

\begin{definition}[Stability on a noncompact hypersurface]\label{def:local-stability}
Let $\Sigma$ be a smooth two-sided minimal hypersurface in a smooth Riemannian manifold, possibly noncompact or with boundary. Put $q=|B|^2+\operatorname{Ric}(\nu,\nu)$ and $J=\Delta_\Sigma+q$. Stability means
\[
 Q(\phi,\phi)=\int_\Sigma(|\nabla\phi|^2-q\phi^2)\,dA\geq0
 \quad\text{for all }\phi\in C_c^\infty(\Sigma\setminus\partial\Sigma).
\]
Thus stability on an open subset only tests variations supported there. The compactly supported second-variation formula is the local version of Theorem \ref{thm:second-variation}; see \cite{ChMin}, Theorem 23.4.
\end{definition}

\begin{proposition}[Linearization of scalar mean curvature]\label{prop:mean-curvature-linearization}
Let $\Sigma$ be smooth, two-sided, and minimal. For a normal variation with velocity $f\nu$, differentiating scalar mean curvature at $t=0$ gives
\[
 \left.\partial_t H_{\Sigma_t}\right|_{t=0}=-Jf.
\]
The sign agrees with $Q(f,f)=-\int fJf$ and with the first variation $\delta A(f\nu)=\int Hf$.
\end{proposition}
\begin{proof}
Take a two-parameter normal graph variation with initial velocities $f\nu$ and $\phi\nu$, both compactly supported. Differentiate its first-variation identity in the other parameter. Terms differentiating the area form or normal component are multiplied by $H=0$, leaving $\int\phi\,\partial_tH$. The mixed second derivative is the polarization of the second variation, $Q(f,\phi)=-\int\phi Jf$. Since $\phi$ is arbitrary, the identity follows pointwise. For a general smooth $f$, insert cutoffs near the point being tested. Compare \cite{ChMin}, the proof of Proposition 33.3, pp.~80--81.
\end{proof}

\begin{proposition}[Positive Jacobi fields imply stability]\label{prop:positive-jacobi-stability}
If $w>0$ is smooth on a two-sided minimal hypersurface and $Jw=0$, then for every compactly supported smooth $\phi$ in its interior,
\[
 Q(\phi,\phi)=\int_\Sigma w^2\left|\nabla\left(\frac\phi w\right)\right|^2dA\geq0.
\]
In particular the hypersurface is stable in the sense of Definition \ref{def:local-stability}.
\end{proposition}
\begin{proof}
Set $\eta=\phi/w$. Expanding $|\nabla(w\eta)|^2$ and subtracting $q w^2\eta^2$, the terms other than $w^2|\nabla\eta|^2$ integrate to
$\int\nabla w\cdot\nabla(w\eta^2)-qw^2\eta^2=-\int w\eta^2(\Delta w+qw)=0$.
There is no boundary term because $\eta$ has compact support. This identity makes explicit the positive-Jacobi-field argument of \cite{ChMin}, Proposition 33.4.
\end{proof}

\begin{corollary}[The vertical angle makes a minimal graph stable]\label{cor:graph-positive-jacobi}
For a smooth minimal graph in $\mathbb R^{n+1}$ with upward normal $\nu$, the function
\[
 w=\langle\nu,e_{n+1}\rangle=(1+|Du|^2)^{-1/2}
\]
is a positive Jacobi function. Hence the graph is stable under all compactly supported normal variations.
\end{corollary}
\begin{proof}
Vertical translations preserve minimality. Their normal speed is $w$; their tangential part differentiates the identically zero mean curvature and contributes zero. Proposition \ref{prop:mean-curvature-linearization} therefore gives $Jw=0$. Apply Proposition \ref{prop:positive-jacobi-stability}. This also follows from the minimizing property in Proposition \ref{prop:graph-calibration}.
\end{proof}

\begin{theorem}[Separated sheets produce a positive Jacobi field]\label{thm:separated-sheets-jacobi}
Let $\Sigma$ be a smooth connected two-sided embedded minimal hypersurface in a smooth Riemannian manifold. Let $U_1\Subset U_2\Subset\cdots$ be connected relatively compact open subsets exhausting its interior, and fix $p\in U_1$. Suppose smooth functions $u_{1,j}<u_{2,j}$ on $U_j$ define minimal normal graphs by the normal exponential map, with both functions converging to zero in $C^\infty$ on every compact subset of $\Sigma\setminus\partial\Sigma$. Then a subsequence of
\[
 v_j=\frac{u_{2,j}-u_{1,j}}{(u_{2,j}-u_{1,j})(p)}
\]
converges in $C^\infty$ on compact subsets to a function $v>0$ with $v(p)=1$ and $Jv=0$. Consequently the interior of $\Sigma$ is stable. This is the normal-graph formulation of \cite{ChMin}, Proposition 33.3, including its stated Riemannian extension. Uniform local ellipticity, Harnack, Schauder estimates, and a diagonal subsequence are the analytic proof inputs; that convergence proof remains external. The conclusion assumes two globally ordered sheets on the exhaustion, not merely a numerical multiplicity at one point.
\end{theorem}

\begin{theorem}[Stable curvature estimate in dimensions at most six]\label{thm:stable-curvature-bound}
Let $2\leq n\leq6$. Suppose $\Sigma^n$ is a smooth properly embedded two-sided minimal hypersurface without boundary in $B_{2R}^{n+1}(a)$, is stable for compactly supported normal variations there, and satisfies
\[
 \mathcal H^n(\Sigma\cap B_{2R}(a))\leq\Lambda R^n.
\]
Then $\sup_{\Sigma\cap B_R(a)}|B|\leq C(n,\Lambda)/R$. This is the smooth Euclidean specialization of Schoen--Simon \cite{SS81}, as stated in \cite{Sharp17}, Theorem 2.1, pp.~320--321. Applying that local estimate at points of $\Sigma\cap B_R(a)$ gives the stated interior bound. Smoothness removes the singular-set hypothesis in that source; for singular hypersurfaces the hypothesis must be restored. The area bound and dimension restriction are retained. The proof remains external.
\end{theorem}
